\documentclass[hyperref={pdfpagelabels=false},table]{beamer}

\input{lec_common}
\begin{document}

\part{Unit 4: Functions}

%\frame{\tableofcontents}


\begin{frame}[fragile]{Functions in Mathematics}
In Mathematics:\\
A function is written as a map, that uniquely\\
assigns an element $y$ from the range $R$
to every element $x$ from the domain $D$.\\
\vfill
\[
 f: x \mapsto y .
\]
$f$ is the function\\
$x$ is its argument\\
$y$ is its (return) value\\
\vfill
There can be several arguments of different type:
\vfill
Consider
\[
f(g,a,b) = \int_a^b g(x) \mathrm{d} x
\]
The arguments are not interchangable. Position matters.
\end{frame}
\begin{frame}[fragile]{Functions in Python}
Definition of a function:
\begin{python}
def f(x1, x2, x3):
    .....  # some computations
    y = ... #
    return y
\end{python}
Evaluation (call) of a function:
\begin{python}
f(17, 18, -2)
f([1, 2, 3], {'Tol': 1.e-10}, 'ro')
...
\end{python}
Wording:\\
\pyth!x1, x2, x3! are called function parameters {\footnotesize (needed for the definition)}\\
\pyth!17, 18, -2! are called arguments {\footnotesize (needed for the evaluation)}
\end{frame}
\begin{frame}[fragile]{Passing Arguments}
Consider:
\begin{python}
def subtract(x1, x2):
    return x1 - x2
\end{python}
\vfill
Passing arguments by position:
\begin{python}
subtract(1, 2) # returns -1
\end{python}
Position matters.
\vfill
Passing arguments by keyword:
\begin{python}
subtract(x2=2, x1=1) # returns -1
\end{python}
Position doesn't matter.
\end{frame}
\begin{frame}[fragile]{Scope of variables I}
Variables defined inside the function
are said to belong to the function's \emph{scope}.\\

They are unknown outside the function.
\vfill
\begin{minipage}[top]{0.3\textwidth}
Definition:
\begin{python}
def mult2(x):
    c = 2.
    return c*x
\end{python}
\end{minipage}
\hfill
\begin{minipage}[top]{0.6\textwidth}
Evaluation:
\begin{python}
mult2(20.)  # returns 40
c    # returns NameError:
     # 'c' not defined
\end{python}
\end{minipage}
\end{frame}
\begin{frame}[fragile]{Scope of variables II}
Compare:\\
\vfill
\begin{minipage}[top]{0.4\textwidth}
\pyth!a! is a parameter of the function
\begin{python}
def multiply(x, a):
    return a*x
\end{python}
\end{minipage}
\hfill
\begin{minipage}[top]{0.4\textwidth}
\pyth!a! is referenced from outside the function's \emph{scope}:
\begin{python}
a = 3
def multiply(x):
    return a*x
\end{python}
\end{minipage}
\end{frame}
\begin{frame}[fragile]{Local variables}
Similarly, immutable variables are copied when calling a function.\\
You cannot change an input variable by changing a local variable like this:
\begin{python}
a = 3
def f(x):
    x = 2
    return x
print(f(a)) # 2
print(a)    # 3
\end{python}
\emph{However}, this \emph{can} be done for mutable variables like lists, dictionaries, etc.\\
As a principle, \emph{never} change input variables within a function.
\end{frame}
\begin{frame}[fragile]{Default Arguments}
{\footnotesize default (eng) = standard (sv)}\ \\
\vfill
\begin{minipage}[top]{0.4\textwidth}
The standard way ...
\begin{python}
def mult(x, c):
    return c*x
\end{python}
a typical call
\begin{python}
y=mult(2., 13.)
# returns 26
y=mult(c=13., x=2.)
# returns 26
\end{python}

\end{minipage}\hfill
\begin{minipage}[top]{0.4\textwidth}
... now with defaults
\begin{python}
def mult(x, c=1.0):
    return c*x
\end{python}
possible calls
\begin{python}
y=mult(2., 13.)
# returns 26
y=mult(2.)
# returns 2
\end{python}
\end{minipage}\\
\vfill
In the definition of the function mandatory parameters \emph{must}
\emph{precede} optional parameters (those with default values).\\
Why?
\end{frame}
\begin{frame}[fragile]{Arguments, Parameters -- Summary}
Make sure that you understood the difference
\begin{itemize}
\item between arguments and parameters
\item between function definition and function evaluation (call)
\item between positional arguments and keyword arguments
\end{itemize}
\end{frame}

\begin{frame}[fragile]{Docstrings}
All functions (and everything else) should be documented
carefully:\\
\vfill
A docstring is the leading comment in a function (or class):\\ \vfill
\begin{python}
def newton(f, x0):
    """
    Newton's method for computing a zero of a function
    on input:
    f  (function) given function f(x)
    x0 (float) initial guess
    on return:
    y  (float) the approximated zero of f
    """
    ...
\end{python}
\pyth!help(newton)! in Python or \pyth!newton?! in IPython
displays the docstring as a help text.
\end{frame}
\begin{frame}[fragile]{Unpacking Arguments }
\emph{Positional arguments} remind us of \emph{lists} \\
\emph{Keyword arguments} remind us of \emph{dictionaries}\\
\vfill
\begin{python}
data = [[1,2], [3,4]]
style = {'linewidth':3, 'marker':'o', 'color':'green'}
\end{python}
\vfill
Star operators unpack these to form a valid parameter list
\begin{python}
plot(*data, **style)
\end{python}
\vfill
$\ast$ unpacks a list to positional arguments\\
$\ast\ast$ unpacks dictionaries to keyword arguments
\end{frame}
\begin{frame}[fragile]{Passing (tunneling) arguments}
 Also in the definition of functions you might find these constructs.
This is often used to pass arguments through a function
\begin{python}
def outer(f, x, *args, **keywords):
   return f(x, *args, **keywords)

def inner(x, y, z, u):
   print(y, z)
   print(u)
   return x**2
\end{python}
\begin{minipage}[top]{.45\textwidth}
 A call
\begin{python}
L=[1, 2]
D={'u':15}
outer(inner, 3, *L, **D)
\end{python}
\end{minipage} \hfill
\begin{minipage}[top]{.5\textwidth}
Equivalently:
\begin{python}
outer(inner, 3, 1, 2, u=15)
\end{python}
\end{minipage}\\
\vfill

Note, the function \pyth!outer! cannot know how many arguments it needs
to provide a full parameter list to the ``inner'' function \pyth!f!.
\end{frame}
\begin{frame}[fragile]{Return}
The \pyth!return! statement returns \emph{a single} object!\\
\vfill
\begin{python}
def my_func(x):
    return 1,2,3,4,5,6
\end{python}
What is the object that is returned here? Which type does it have? \\
{\footnotesize (see Unit 3)}\\
\vfill
Statements after the return statement are ignored:\\[0.5cm]
\begin{minipage}[top]{0.4\textwidth}
\begin{python}
def my_func(x):
    return 1,2,3
    z = 25 # ignored
\end{python}
\end{minipage}
\hfill
\begin{minipage}[top]{0.4\textwidth}
\begin{python}
def my_func(x):
   if x > 0:
      return 1
   else:
      return -1
   z = 25 # ignored
\end{python}
\end{minipage}
\end{frame}
\begin{frame}[fragile]{No Return}
A function \emph{without} a \pyth!return! statement
returns \pyth!None!:\\
\vfill
\begin{python}
def my_func(x):
    z = 2*x

a = my_func(10.)

type(a) # <type 'NoneType'>
a == None # true
\end{python}
\vfill
\end{frame}
\begin{frame}[fragile]{Functions are objects}
Functions are objects, they can be deleted, reassigned, copied ... \\
\vfill
\begin{python}
def square(x):
  """Return the square of `x`"""
  return x**2
square(4) # 16
sq = square # now sq is the same as square
sq(4) # 16
del square # `square` doesn't exist anymore
print(newton(sq, .2)) # passing as argument
\end{python}
\end{frame}
\begin{frame}[fragile]{Partial Application}
{\footnotesize Partial application = closures}\ \\
\vfill
In mathematics we often ``freeze'' a parameter of a function:
\[
 f_{\omega}(x)=f(x,\omega)=\sin(\omega x)
\]
\vfill
In Python there are many possibilities to do this, here is one ...
\begin{python}
def make_sin(omega):
    def f(x):
        return sin(omega*x)
    return f

fomega = make_sin(omega)

\end{python}
In order to understand this solution, recall
\begin{itemize}
\item the scope of a variable and references to variables out of scope
\item functions are objects
\end{itemize}
\end{frame}
\begin{frame}[fragile]{Anonymous Functions: the \pyth{lambda} keyword}
With $\lambda$-functions one has a handy tool for
making one-line function definitions:
\begin{python}
f = lambda x: 3.*x**2 + 2.*x + 0.5
f(3)  # returns 33.5
g = lambda x,y: 3.*x-2.*y
g(1,1) # returns 1.0
\end{python}
\vfill
Example for a common application, compute $\int_0^1 x^2 + 5 \, \mathrm{d} x$:\\
\vfill
\begin{python}
import scipy.integrate as si
si.quad(lambda x: x**2+5,0,1)
\end{python}
\end{frame}
\begin{frame}[fragile]{Partial Applications and $\lambda$}
\[
 f_{\omega}(x)=f(x,\omega)=\sin(\omega x)
\]
\vfill
 .... now simply becomes\\
\vfill
\begin{python}
omega = 3.
fomega = lambda x: sin(omega*x)
fomega(1.) # returns 0.14112...
\end{python}

\end{frame}

\end{document}
